\documentclass[9pt,a4paper]{extarticle}

% ======================
% Pacotes básicos
% ======================
\usepackage[T1]{fontenc}
\usepackage[utf8]{inputenc}
\usepackage[brazil]{babel}
\usepackage{geometry}
\geometry{margin=2cm}
\usepackage{parskip}
\usepackage{setspace}
\usepackage{microtype}

% ======================
% Matemática
% ======================
\usepackage{amsmath, amssymb, amsthm}
\usepackage{bm}

% ======================
% Listas e hyperlinks
% ======================
\usepackage{enumitem}
\setlist{nosep}
\usepackage{hyperref}
\hypersetup{colorlinks=true,linkcolor=black,urlcolor=blue}

% ======================
% Título
% ======================
\title{\vspace{-1cm}Roteiro de Aula -- Métodos de Gradiente para Treinamento de Redes Neurais}
\author{Prof.\ Rodrigo Silva}
\date{}

\begin{document}
\maketitle

\section{Motivação}
O treinamento de redes neurais consiste em minimizar uma função de perda 
\[
L(\theta),
\]
onde $\theta$ representa todos os parâmetros do modelo. O gradiente pode ser calculado eficientemente via \emph{backpropagation}, de modo que métodos baseados em gradiente são a principal ferramenta de otimização em aprendizagem profunda.

\vspace{0.2cm}
Nesta aula estudaremos os principais métodos de descida comumente utilizados no treinamento de redes neurais.

% ============================================================
\section{Gradient Descent Básico}

Dado um vetor de parâmetros $\theta_k$, a atualização por descida do gradiente é:

\[
\theta_{k+1} = \theta_k - \alpha \nabla_\theta L(\theta_k),
\]

onde $\alpha > 0$ é a taxa de aprendizado. As principais variantes práticas são:

\begin{itemize}
    \item \textbf{Batch GD}: usa todo o conjunto de dados.
    \item \textbf{Stochastic GD}: usa apenas um exemplo por iteração.
    \item \textbf{Mini-batch GD}: usa subconjuntos dos dados (forma usual).
\end{itemize}

Limitações principais:
\begin{itemize}
    \item sensibilidade à escolha do $\alpha$;
    \item pode oscilar em vales estreitos da função;
    \item pode convergir lentamente.
\end{itemize}

% ============================================================
\section{Gradient Descent com Momentum}

O método adiciona o conceito de ``velocidade'' acumulada:

\[
v_{k+1} = \beta v_k + (1-\beta)\nabla L(\theta_k),
\]
\[
\theta_{k+1} = \theta_k - \alpha v_{k+1},
\]

onde $\beta \in [0.9,0.99]$ regula a influência do passado.

Vantagens:
\begin{itemize}
    \item acelera o avanço em direções consistentes;
    \item reduz oscilações devido a ruído no gradiente.
\end{itemize}

% ============================================================
\section{AdaGrad}

O AdaGrad adapta a taxa de aprendizado individualmente para cada parâmetro acumulando os quadrados dos gradientes:

\[
g_{k+1} = g_k + \left(\nabla L(\theta_k)\right)^2,
\]

\[
\theta_{k+1} = \theta_k - \frac{\alpha}{\sqrt{g_{k+1}} + \varepsilon} \nabla L(\theta_k).
\]

Consequências:
\begin{itemize}
    \item parâmetros com gradientes grandes têm passo reduzido;
    \item parâmetros com gradientes pequenos têm passo ampliado.
\end{itemize}

Problema principal: o termo acumulado cresce monotonicamente, levando a passos cada vez menores.

% ============================================================
\section{RMSProp}

RMSProp corrige o AdaGrad substituindo a soma acumulada por uma média móvel exponencial:

\[
s_{k+1} = \beta s_k + (1-\beta)(\nabla L(\theta_k))^2,
\]

\[
\theta_{k+1} = \theta_k - \frac{\alpha}{\sqrt{s_{k+1}} + \varepsilon} \nabla L(\theta_k).
\]

Vantagens:
\begin{itemize}
    \item evita que o passo tenda a zero;
    \item funciona bem com gradientes ruidosos;
    \item é amplamente utilizado em redes profundas.
\end{itemize}

% ============================================================
\section{RMSProp com Momentum}

Combinação da média móvel adaptativa do RMSProp com o acúmulo de velocidade do momentum:

\[
s_{k+1} = \beta s_k + (1-\beta)(\nabla L(\theta_k))^2,
\]
\[
v_{k+1} = \mu v_k + \frac{\alpha}{\sqrt{s_{k+1}}+\varepsilon}\nabla L(\theta_k),
\]
\[
\theta_{k+1} = \theta_k - v_{k+1}.
\]

Vantagens: maior estabilidade e aceleração.

% ============================================================
\section{AdaDelta}

O AdaDelta busca eliminar a dependência direta do parâmetro $\alpha$.

Acumula gradientes:
\[
E[g^2]_{k+1} = \beta E[g^2]_k + (1-\beta)(\nabla L(\theta_k))^2.
\]

Define a atualização:
\[
\Delta\theta_k = -\frac{\sqrt{E[\Delta\theta^2]_k + \varepsilon}}{\sqrt{E[g^2]_{k+1}}+\varepsilon}\nabla L(\theta_k).
\]

Acumula as atualizações:
\[
E[\Delta\theta^2]_{k+1} = \beta E[\Delta\theta^2]_k + (1-\beta)(\Delta\theta_k)^2.
\]

Benefícios:
\begin{itemize}
    \item reduz a necessidade de escolher $\alpha$;
    \item acelera regiões planas e estabiliza regiões de grande curvatura.
\end{itemize}

% ============================================================
\section{Adam -- Adaptive Moment Estimation}

O Adam combina momentum com RMSProp e inclui correção de viés.

Momentos:
\[
m_{k+1} = \beta_1 m_k + (1-\beta_1)\nabla L(\theta_k),
\]
\[
v_{k+1} = \beta_2 v_k + (1-\beta_2)(\nabla L(\theta_k))^2.
\]

Correções de viés:
\[
\hat{m}_{k+1} = \frac{m_{k+1}}{1-\beta_1^{k+1}}, \qquad
\hat{v}_{k+1} = \frac{v_{k+1}}{1-\beta_2^{k+1}}.
\]

Atualização final:
\[
\theta_{k+1} = \theta_k - \frac{\alpha \hat{m}_{k+1}}{\sqrt{\hat{v}_{k+1}} + \varepsilon}.
\]

Vantagens principais:
\begin{itemize}
    \item rápido e estável;
    \item altamente utilizado em aprendizagem profunda;
    \item adapta-se automaticamente a diferentes escalas de parâmetros.
\end{itemize}

% ============================================================
\section{Resumo Comparativo}

\begin{itemize}
    \item GD: simples, porém lento e sensível ao passo.
    \item Momentum: acelera e suaviza oscilações.
    \item AdaGrad: adapta o passo, mas pode estagnar.
    \item RMSProp: adapta o passo sem estagnação.
    \item RMSProp + Momentum: equilíbrio entre estabilidade e velocidade.
    \item AdaDelta: reduz dependência do $\alpha$.
    \item Adam: método mais utilizado atualmente.
\end{itemize}

\end{document}
